% Generated draft: requires mathematical and source review.
\documentclass[11pt]{article}
\usepackage[margin=25mm]{geometry}
\usepackage{fontspec}
\setmainfont{Latin Modern Roman}
\usepackage{amsmath,amssymb,mathtools,unicode-math}
\setmathfont{Latin Modern Math}
\usepackage{enumitem,needspace,xurl,hyperref}
\setlength{\parindent}{0pt}
\setlength{\parskip}{5pt}
\setlength{\emergencystretch}{3em}
\newcommand{\N}{\mathbb{N}}
\newcommand{\Z}{\mathbb{Z}}
\newcommand{\Q}{\mathbb{Q}}
\newcommand{\R}{\mathbb{R}}
\newcommand{\C}{\mathbb{C}}
\newcommand{\F}{\mathbb{F}}
\newcommand{\A}{\mathbb{A}}
\newcommand{\PP}{\mathbb{P}}
\newcommand{\E}{\mathbb{E}}
\newcommand{\eps}{\varepsilon}
\newcommand{\dd}{\,\mathrm{d}}
\begin{document}
\section*{UnsolvedMath ID: 30004798}
\textbf{Source problem code:} catalogue entry
\subsection{Definitions and mathematical statement}
Let \(F\) be a number field and let \(\mathbb A_F\) denote its ring of adeles. Let \(G\) be an even orthogonal or unitary group over \(F\). For automorphic representations of \(G(\mathbb A_F)\), near-equivalence means that their local components are isomorphic at all but finitely many places of \(F\).

Let \(\psi\) be an elliptic Arthur parameter, also called an elliptic \(A\)-parameter, of \(G\). Write \(\mathcal N_\psi\) for the near-equivalence class represented by \(\psi\), as specified in the supplied statement. The record does not give the definition of these parameters or the construction assigning this class to \(\psi\).

For \(\pi\in\mathcal N_\psi\), retain the source notation \(m_{\mathrm{disc}}(\pi)\) and \(m_{\mathrm{aut}}(\pi)\) for its discrete and automorphic multiplicities, respectively. The record does not define the representation spaces or conventions used to compute these multiplicities; in particular, it does not provide a mathematical definition of \(m_{\mathrm{aut}}\).

\textbf{Conjecture as supplied.} For every number field \(F\), every even orthogonal or unitary group \(G\) over \(F\), every elliptic Arthur parameter \(\psi\) of \(G\), and every automorphic representation \(\pi\) in the near-equivalence class represented by \(\psi\),
\[
 m_{\mathrm{disc}}(\pi)=m_{\mathrm{aut}}(\pi).
\]

\textbf{Unresolved qualifications.} This is a transcription of the supplied conjecture, conditional on the source's intended definitions. The record does not specify whether the orthogonal group is the full orthogonal group or the special orthogonal group, any restrictions on its forms, the quadratic extension and Hermitian data for the unitary case, or any central-character conventions. These choices, the precise parameter formalism, and the definitions of both multiplicities require clarification before this can be treated as a fully specified standalone statement. No additional restrictions are imposed here.
\subsection{Short English statement}
The supplied conjecture asserts that every automorphic representation in the near-equivalence class associated with an elliptic Arthur parameter of an even orthogonal or unitary group has equal discrete and automorphic multiplicities. The meanings of the multiplicities and several group and parameter conventions remain unspecified. The supplied dated curation labels the problem partially solved; this draft does not independently verify that assessment.
\subsection{Sources}
\begin{itemize}
\item Local registry: \texttt{data/unsolvedmath/problems.json}, permanent UnsolvedMath ID \(30004798\), title \emph{Equality of Discrete and Automorphic Multiplicities}, source problem code \texttt{catalogue entry}, in the Oberwolfach Reports Open Problems collection. The supplied record lists proposed year 2021 and names no proposer.
\item Supplied report citation: \emph{Automorphic Forms, Geometry and Arithmetic (2022)}, Oberwolfach Reports, pp.~2089--2156, DOI \texttt{10.4171/owr/2021/39}. \url{https://doi.org/10.4171/owr/2021/39}. The background describes Oberwolfach Reports as documenting mathematical workshops and their open-problem discussions. The citation's year 2022 and the record's proposed year 2021 are retained as supplied.
\item Supplied literature reference: Chen--Zou, \emph{Arthur multiplicity formula}. \url{https://ems.press/journals/jems/articles/14298407}.
\item Supplied literature reference: Chen--Zou, \emph{theta correspondence and Arthur packets}. \url{https://arxiv.org/abs/2104.12354}.
\item Supplied literature assessment, checked at \texttt{2026-08-22}: status \texttt{partially\_solved}. The curation attributes to Chen and Zou a proof of the multiplicity formula in the stated broad families of even orthogonal and unitary groups, including the low-Witt-index case, while saying that the extracted conjecture is not resolved for every group and parameter. The record does not spell out those families or the low-Witt-index hypotheses. This is supplied dated curation, not an independently verified literature assessment; its time-sensitive status and cited claims require independent verification.
\item Supplied statement-verification metadata: status \texttt{corrected\_verified}, with the assertion that the standalone statement and its hypotheses were checked against the official Oberwolfach report. This is a report of the supplied metadata, not a verification performed for this draft.
\end{itemize}
\end{document}
